\section{三个典型应用域}

\subsection{Coding Agent：代码库是可执行环境}

编码任务具有天然反馈：编译器、测试、静态分析和运行日志都能充当 verifier。Coding agent 的价值不仅是补全代码，而是维护一个持续变化的工作状态：搜索仓库、编辑多个文件、运行测试、读取失败、定位根因并迭代。

\begin{figure}[H]
\centering
\includegraphics[width=0.94\textwidth]{figures/coding-agents.jpg}
\caption{Coding agent 同时操作 shell、浏览器、编辑器和计划状态。\protect\footnotemark}
\end{figure}
\footnotetext{视频讲解区间：00:52:44--00:54:34。}

\begin{knowledgebox}{为什么代码是自我改进方法的理想试验场}
代码任务同时具备开放式生成与相对便宜的自动验证：候选程序空间巨大，但单元测试、类型检查和执行结果可快速过滤候选。这使代码成为 repeated sampling、search、execution feedback 和 RL 的高价值交汇点。
\end{knowledgebox}

\subsection{Customer Support Agent：反馈来自业务状态}

客服 Agent 需要访问订单、工单、退款和知识库，并执行真实业务动作。成功标准通常不是“回答听起来合理”，而是问题是否解决、状态是否正确更新、用户是否需要再次联系。

\begin{figure}[H]
\centering
\includegraphics[width=0.93\textwidth]{figures/customer-support-agents.jpg}
\caption{客服 Agent 把信息检索、判断和业务动作连接起来。\protect\footnotemark}
\end{figure}
\footnotetext{视频讲解区间：00:54:34--00:55:43。}

这类系统的难点是权限与不可逆动作：模型可以建议退款，但真正执行前往往需要规则检查、额度限制或人工确认。

\subsection{Research Agent：把研究过程分成可验证阶段}

科研代理覆盖想法生成、文献检索、实验设计、代码执行、结果分析和论文写作。不同阶段的 verifier 差异很大：检索可检查引用是否存在，代码可运行测试，实验结论却需要统计严谨性和领域判断。

\begin{figure}[H]
\centering
\includegraphics[width=0.96\textwidth]{figures/research-agents.jpg}
\caption{研究 Agent 的想法、实验迭代和论文写作流程。\protect\footnotemark}
\end{figure}
\footnotetext{视频讲解区间：00:56:00--00:59:03。}

\begin{warningbox}{“自动完成报告”不等于完成研究}
Research agent 最容易在引用真实性、实验可复现性和因果解释上产生表面完整但证据薄弱的结果。高质量系统必须把来源核验、实验日志、统计检验和反例搜索纳入工作流，而不能只优化最终文本的流畅度。
\end{warningbox}

\subsection{本章小结}

三类应用的共同点是：模型必须离开纯文本空间，与可变化的外部状态交互。它们的差异主要来自 verifier 的成本与可靠性：代码执行反馈最便宜，业务成功信号更延迟，科研正确性最难自动化判断。

